% ============================================================
%  KATA PENGANTAR — UI Skripsi (3 penulis)
%  Compile standalone:  pdflatex kataPengantar.tex
%  Compile as part of thesis: pdflatex thesis.tex
%
%  Diatur sebagai \subfile dari thesis.tex sehingga preamble (font,
%  geometri, setspace, hyperref) otomatis terwarisi. Bila perlu kompilasi
%  mandiri, kelas \texttt{subfiles} akan menggunakan preamble thesis.tex.
%
%  TODO (revisi final):
%    * Isi nama Pembimbing Akademik Fadrian Yhoga Pratama.
%    * Lengkapi nama \emph{orang terkasih} masing-masing penulis (butir 4).
%    * Pertimbangkan apakah nama enam pakar validasi disebut eksplisit
%      atau tetap anonim selaras dengan Lampiran 8.
% ============================================================
\documentclass[../../thesis]{subfiles}

\begin{document}

\addcontentsline{toc}{chapter}{KATA PENGANTAR}
\begin{center}\bfseries KATA PENGANTAR\end{center}
\vspace{0.8cm}
\setstretch{1.4}

\noindent Puji syukur kami panjatkan kepada Tuhan Yang Maha Esa, karena atas
berkat, rahmat, dan karunia-Nya, kami dapat menyelesaikan skripsi ini.
Penulisan skripsi ini dilakukan dalam rangka memenuhi salah satu syarat untuk
mencapai gelar Sarjana Ilmu Komputer pada Fakultas Ilmu Komputer, Universitas
Indonesia. Kami menyadari bahwa, tanpa bantuan dan bimbingan dari berbagai
pihak, sejak masa perkuliahan hingga penyusunan skripsi ini, sangatlah sulit
bagi kami untuk menyelesaikannya. Oleh karena itu, kami mengucapkan terima
kasih kepada:

\begin{enumerate}
  \item Tuhan Yang Maha Esa, atas berkat, rahmat, kesehatan, dan kemudahan
        yang tiada henti dilimpahkan kepada kami selama proses penyusunan
        skripsi ini;
  \item Dr.\ Siti Aminah, S.Kom., M.Kom., selaku dosen pembimbing skripsi,
        yang dengan sabar telah menyediakan waktu, tenaga, dan pikiran untuk
        mengarahkan kami sejak tahap perumusan masalah hingga penyusunan
        akhir naskah skripsi ini;
  \item Dosen Pembimbing Akademik kami masing-masing yang telah membimbing
        perjalanan akademik kami sepanjang masa perkuliahan, yaitu:
        \begin{enumerate}
          \item \emph{[Nama Pembimbing Akademik --- diisi pada revisi final]},
                selaku Pembimbing Akademik Fadrian Yhoga Pratama;
          \item Dr.\ Ade Azurat, S.Kom., selaku Pembimbing Akademik
                Soros Febriano; serta
          \item Dr.\ Siti Aminah, S.Kom., M.Kom.\ --- yang juga merupakan
                dosen pembimbing skripsi kami --- selaku Pembimbing Akademik
                Tegar Wahyu Khisbulloh;
        \end{enumerate}
  \item Orang-orang terkasih yang selalu setia mendampingi, menyemangati,
        dan menjadi tempat berbagi suka maupun duka kami masing-masing
        sepanjang proses penyusunan skripsi
        ini\footnote{Nama-nama secara spesifik dicantumkan
        pada revisi final masing-masing penulis.};
  \item Para pakar bidang Fisika, Kimia, dan Biologi yang telah bersedia
        meluangkan waktu untuk berperan sebagai validator domain pada proses
        evaluasi pakar dalam penelitian ini --- waktu, ketelitian, dan
        kepakaran mereka sangat berarti bagi kualitas hasil penelitian ini;
  \item Orang tua dan keluarga besar kami masing-masing, yang telah
        memberikan dukungan material, moral, dan doa yang tiada henti sejak
        awal masa studi hingga selesainya skripsi ini;
  \item Teman-teman seangkatan, sahabat, serta semua pihak yang tidak dapat
        kami sebutkan satu per satu, yang telah membantu kami baik secara
        langsung maupun tidak langsung selama masa studi dan penyusunan
        skripsi ini; serta
  \item Sesama anggota tim peneliti --- Fadrian Yhoga Pratama, Soros Febriano,
        dan Tegar Wahyu Khisbulloh --- atas kerja sama, dedikasi, kesabaran,
        dan kepercayaan yang terjaga sepanjang perjalanan penelitian ini.
\end{enumerate}

\noindent Akhir kata, kami berharap Tuhan Yang Maha Esa berkenan membalas
segala kebaikan semua pihak yang telah membantu. Semoga skripsi ini membawa
manfaat bagi pengembangan ilmu pengetahuan, khususnya bagi upaya pemetaan
interkoneksi materi sains pada Kurikulum Merdeka.
\vspace{1cm}

\begin{flushright}
Depok, Juni 2026\\[1.5cm]
Penulis
\end{flushright}

\end{document}
